% Generated draft: requires mathematical and source review.
\documentclass[11pt]{article}
\usepackage[margin=25mm]{geometry}
\usepackage{fontspec}
\setmainfont{Latin Modern Roman}
\usepackage{amsmath,amssymb,mathtools,unicode-math}
\setmathfont{Latin Modern Math}
\usepackage{enumitem,needspace,xurl,hyperref}
\setlength{\parindent}{0pt}
\setlength{\parskip}{5pt}
\setlength{\emergencystretch}{3em}
\newcommand{\N}{\mathbb{N}}
\newcommand{\Z}{\mathbb{Z}}
\newcommand{\Q}{\mathbb{Q}}
\newcommand{\R}{\mathbb{R}}
\newcommand{\C}{\mathbb{C}}
\newcommand{\F}{\mathbb{F}}
\newcommand{\A}{\mathbb{A}}
\newcommand{\PP}{\mathbb{P}}
\newcommand{\E}{\mathbb{E}}
\newcommand{\eps}{\varepsilon}
\newcommand{\dd}{\,\mathrm{d}}
\begin{document}
\section*{UnsolvedMath ID: 30001901}
\textbf{Source problem code:} catalogue entry
\subsection{Definitions and mathematical statement}
A number field is a field \(K\) that is a finite extension of \(\Q\). Two number fields \(K\) and \(L\) are isomorphic if there exists a field isomorphism \(K\to L\); such an isomorphism necessarily fixes \(\Q\).

The supplied question asks whether an invariant internal to a number field can classify number fields up to isomorphism and provide an internal analogue of the external reconstruction results mentioned in the source record. Schematically, after specifying a class of invariant values \(\mathcal X\), an equivalence relation \(\sim\) on those values, and the admissible meaning of ``internal,'' the intended classification requirement would be an assignment \(K\mapsto I(K)\in\mathcal X\) satisfying
\[
\forall\,\text{number fields }K,L,\qquad
I(K)\sim I(L)\quad\Longleftrightarrow\quad K\cong L.
\]
The assignment would additionally have to satisfy the intended internality condition and the intended relationship to the external reconstruction results.

\textbf{Unresolved formulation.} The record does not define internality, the permitted invariant values or their equivalence, or the external reconstruction results to be matched. These are unresolved parts of the question, not freely chosen hypotheses. The displayed condition therefore expresses only the classification requirement, not a complete formalization of the intended problem. An unrestricted complete invariant would not by itself resolve the unspecified internality requirement.
\subsection{Short English statement}
The question asks whether data internal to each number field can determine that field up to isomorphism, in a way analogous to unspecified external reconstruction results. The supplied fragment does not define the permitted internal data or the required analogy, so the intended existence question remains incompletely specified.
\subsection{Sources}
\begin{itemize}
\item Supplied local UnsolvedMath registry record, permanent ID \(30001901\), title: \emph{Internal Isomorphism Invariants of Number Fields}; source problem code: \texttt{catalogue entry}. The record belongs to the Oberwolfach Reports Open Problems collection and lists proposed year \(2011\); no individual proposer is supplied.
\item Supplied report citation: \emph{Noncommutative Geometry (2012)}, Oberwolfach Reports, DOI \texttt{10.4171/owr/2011/45}, pp.\ 2549--2622. Source URL and supplied literature source, labelled ``original OWR report'': \url{https://doi.org/10.4171/owr/2011/45}. The background describes Oberwolfach Reports as documenting mathematical workshops and their open-problem discussions. No more specific passage or reconstruction theorem is identified.
\item Supplied literature assessment, checked on \texttt{2026-08-21}: the status is labelled open; the extracted sentence is described as a context fragment that does not state a precise resolved theorem, and the curation reports that no authoritative source was located giving a definitive answer to the implied question. The listed source checked is the original OWR report above. This is supplied dated curation, not an independently verified assessment; the record requests independent verification of time-sensitive status and cited claims.
\item The record labels the statement ``exact'' and reports that its mathematical content was checked against the official report, with notation, context, and question form normalized without changing the source-level claim. This is a supplied verification claim, not verification performed for this draft.
\end{itemize}
\end{document}
